\documentclass[10pt,a4paper]{amsart}
\usepackage[margin=19mm]{geometry}
\usepackage{amsmath,amssymb,mathtools}
\usepackage[scr=rsfs,cal=euler]{mathalfa}
\usepackage{xfrac}
\usepackage[unicode,hidelinks,bookmarks=false]{hyperref}
\newcommand{\ie}{i.e.\ }
\newcommand{\cf}{cf.\ }
\newcommand{\resp}{respectively\ }
\newcommand{\Subsection}{\subsection}
\providecommand{\nobreakdash}{\nobreak\hbox{-}}
\setlength{\emergencystretch}{2em}
\multlinegap0pt
\usepackage{amssymb}
\usepackage{mathtools}
\usepackage[shortlabels]{enumitem}
\usepackage{xcolor}
\usepackage[normalem]{ulem}
\usepackage{float}
\newtheorem{lemma}{Lemma}
\newtheorem{theorem}{Theorem}
\title{Two-point completion of zero interlacing and Wronskian-type bounds, with applications to Jacobi and other orthogonal polynomials}
\author{Kerstin Jordaan, Vikash Kumar}
\date{Source: arXiv:2604.25692v2 (CC BY 4.0)}
\begin{document}
\maketitle

\noindent\emph{Source and licence.} Two-point completion of zero interlacing and Wronskian-type bounds, with applications to Jacobi and other orthogonal polynomials. Version arXiv:2604.25692v2 (CC BY 4.0), distributed by arXiv under the Creative Commons Attribution 4.0 International licence, http://creativecommons.org/licenses/by/4.0/, as recorded in the arXiv licence metadata for this exact version and shown on the abstract page https://arxiv.org/abs/2604.25692v2. Full licence notice: \texttt{licenses/arxiv-2604.25692-CC-BY-4.0.txt}.\par\smallskip

\noindent\textit{Editorial scope.} Counted unit: the article's theorem MainTheorem4 (source0.tex lines 234--240). The article's complete original proof of it is delivered with it (source0.tex lines 242--349, one \texttt{proof} environment of 108 lines). No mathematics is written, repaired or re-derived: the statement and the proof are the article's own lines, in the article's own notation, and no retained line is substituted or re-flowed.  Retained uncounted as the source-local context the proof needs: lines 154--154; lines 174--174.  Nothing was omitted from any retained span, so the package carries no omission record.  No retained line contains a citation, so the wrapper carries no bibliography and no bibliography entry.  No retained line contains a figure environment, an image-inclusion command or drawing code, so no image is delivered and the package declares no asset dependency.  Editorial formatting only, in addition: an AMS \texttt{amsart} preamble that restates the article's own declarations and macros used by the retained lines, namely the environments \texttt{lemma}, \texttt{theorem} on separate counters, together with the standard packages those lines need.  The retained environments print in this document as Lemma 1, Theorem 1 in order of appearance, whereas the article prints the same items with the numbering of its own full document.  The complete native source of the article as distributed in the arXiv e-print for this exact version is bundled unchanged as \texttt{sources/199322/source0.tex}.
\begin{lemma}\label{lem:parity-counting} Let $P_n$, $G_{n+1}$, and $Q_n$ be monic polynomials of degrees $n$, $n+1$, and $n$, respectively, whose zeros are real and simple. Assume that $G_{n+1}\prec Q_n$ and that, for all real $x$, \begin{equation} A(x)P_n(x) = B(x)G_{n+1}(x) + \varepsilon (x-E_1)(x-E_2)Q_n(x), \qquad \varepsilon\in\{-1,1\}, \label{eq:general-signed-recurrence} \end{equation} where $A(x)>0$, $E_1,E_2\in \mathbb{R}$, $E_1<E_2$. Suppose further that $P_n$ and $G_{n+1}$ are coprime and that $B(E_i)\neq 0$, $i=1,2.$ Then $F_{n+2}$ has an odd number of zeros in every interior interval $I_k$, $k=1,\dots,n$. Moreover, the number of zeros of $F_{n+2}$ in each exterior interval $I_0$ and $I_{n+1}$ is odd when $\varepsilon=-1$ and even when $\varepsilon =1$.\end{lemma} 

\begin{theorem}\label{MainTheorem3} Let $P_n$, $G_{n+1}$, and $Q_n$ be monic polynomials of degrees $n$, $n+1$, and $n$, respectively, with real and simple zeros. Assume that $G_{n+1}\prec Q_n$ and that, for all real $x$, \begin{equation} A(x)P_n(x) = B(x)G_{n+1}(x) - (x-E_1)(x-E_2)Q_n(x), \qquad n\geq1, \label{GeneralMixedTTRR3} \end{equation} where $E_1,E_2\in\mathbb{R}$, $E_1<E_2$, and $A(x)>0$. Suppose also that $B(E_i)\neq0$, $ i=1,2$. If $P_n$ and $G_{n+1}$ are coprime, then \[ (x-E_1)(x-E_2)P_n(x)\prec G_{n+1}(x). \] If, in addition, both $E_1$ and $E_2$ lie outside the interval determined by the extreme zeros of $G_{n+1}$, then necessarily \[ E_1<y_{1,n+1}, \qquad E_2>y_{n+1,n+1}, \] and \[ G_{n+1}(x)\prec P_n(x). \] If $P_n$ and $G_{n+1}$ are not coprime, then every common zero belongs to $\{E_1,E_2\}$. In particular, $P_n$ and $G_{n+1}$ can have at most two common zeros. \end{theorem} 

\begin{theorem}\label{MainTheorem4} Assume that all the hypotheses of Theorem \ref{MainTheorem3} hold, except that \eqref{GeneralMixedTTRR3} is replaced by \begin{equation} A(x)P_n(x) = B(x)G_{n+1}(x) + (x-E_1)(x-E_2)Q_n(x). \label{GeneralMixedTTRR4} \end{equation} Assume that $P_n$ and $G_{n+1}$ are coprime. Then the following statements hold. \begin{enumerate} \item[(a)] The configuration $ E_1<y_{1,n+1}$ and $E_2>y_{n+1,n+1} $ is impossible. \item[(b)] If $E_1$ and $E_2$ lie in the same one of the intervals $I_0,I_1,\ldots,I_n,I_{n+1}$, then $G_{n+1}(x)\prec P_n(x)$. \item[(c)] If $ E_1<y_{1,n+1} $ and $ y_{k,n+1}<E_2<y_{k+1,n+1} $ for some $k\in\{1,\ldots,n\}$, then $(x-E_2)P_n(x)\prec G_{n+1}(x).$ \item[(d)] If $E_2>y_{n+1,n+1}$ and $y_{k,n+1}<E_1<y_{k+1,n+1}$ for some $k\in\{1,\ldots,n\}$, then \newline $G_{n+1}(x)\prec (x-E_1)P_n(x)$. 
\item[(e)] Let $z_{1,n}<z_{2,n}<\cdots<z_{n,n}$ be the zeros of $P_n$.   Suppose that $E_1$ and $E_2$ lie in two distinct interior intervals \[ y_{k,n+1}<E_1<y_{k+1,n+1}, \qquad y_{k',n+1}<E_2<y_{k'+1,n+1}, \qquad k\neq k'. \] Then, for $l\in\{1,\ldots,n-1\}$, \begin{itemize}\item[(i)]    $(x-E_1)G_{n+1}(x)\prec (x-E_2)P_n(x)$  whenever $
				y_{k,n+1} < z_{l,n} < E_1 < z_{l+1,n} < y_{k+1,n+1}
				$, 
				\item[(ii)] $(x-E_2)G_{n+1}(x)\prec (x-E_1)P_n(x)$  whenever $
				y_{k',n+1} < z_{l,n} < E_2< z_{l+1,n} < y_{k'+1,n+1}$. 
			\end{itemize}\end{enumerate} \end{theorem} 

\begin{proof}
Set
\[
F_{n+2}(x)=(x-E_1)(x-E_2)P_n(x).
\]
Applying Lemma~\ref{lem:parity-counting} with $\varepsilon=1$, we find that $F_{n+2}$
has an odd number of zeros in each interior interval
\[
I_1,\ldots,I_n
\]
and an even number of zeros in each exterior interval $I_0$ and
$I_{n+1}$. In particular, every interior interval contains at least one zero of
$F_{n+2}$. These account for at least $n$ zeros. Since
$F_{n+2}$ has degree $n+2$, there are only two further zeros.
Moreover, adding zeros to any one interval without changing the parity
prescribed by Lemma~\ref{lem:parity-counting} requires adding them in pairs. Therefore, the
two further zeros must lie in the same interval. We call this interval
the exceptional interval. Thus, either one exterior interval contains
exactly two zeros, or one interior interval contains exactly three
zeros. 

\medskip \noindent Suppose first that
\[
E_1<y_{1,n+1}
\qquad\text{and}\qquad
E_2>y_{n+1,n+1}.
\]
Then each exterior interval contains one of the completion points.
Since the number of zeros in each exterior interval must be even, each
would have to contain at least two zeros. This would require at least
four zeros in the exterior intervals, in addition to at least one zero
in each of the $n$ interior intervals, contradicting
$\deg F_{n+2}=n+2$. This proves part~(a).

\medskip \noindent Suppose next that $E_1$ and $E_2$ lie in the same interval. If they
lie in an exterior interval, that interval already contains two zeros
of $F_{n+2}$, so it is the exceptional interval. If they lie in an
interior interval, that interval contains at least the two zeros
$E_1$ and $E_2$; since its number of zeros must be odd, it contains
at least three zeros and is again the exceptional interval. In
either case, after removing the two completion points, $P_n$ has
exactly one zero in every interior interval and no zero in either
exterior interval. Hence $G_{n+1}(x)\prec P_n(x),$ which proves part~(b).



\medskip \noindent Now suppose that
\[
E_1\in I_0
\qquad\text{and}\qquad
E_2\in I_k
\]
for some $k\in\{1,\ldots,n\}$. Since $I_0$ contains $E_1$ and
must contain an even number of zeros of $F_{n+2}$, it contains at
least one further zero. Thus $I_0$ is the exceptional interval.
It follows that $I_0$ contains exactly two zeros, every interior
interval contains exactly one zero, and $I_{n+1}$ contains none. The unique zero of $F_{n+2}$ in $I_k$ is $E_2$, so $P_n$ has no
zero in $I_k$. The other zero of $F_{n+2}$ in $I_0$, besides
$E_1$, is a zero of $P_n$, and every interior interval other than
$I_k$ contains exactly one zero of $P_n$. Therefore
$(x-E_2)P_n(x)$ has exactly one zero in $I_0$ and in every interior
interval, and no zero in $I_{n+1}$. Hence
\[
(x-E_2)P_n(x)\prec G_{n+1}(x),
\]
which proves part~(c).

\medskip \noindent The proof of part~(d) is analogous. If
\[
E_1\in I_k
\qquad\text{and}\qquad
E_2\in I_{n+1},
\]
then $I_{n+1}$ is the exceptional interval. Thus
$(x-E_1)P_n(x)$ has no zero in $I_0$, exactly one zero in every
interior interval, and exactly one zero in $I_{n+1}$. Consequently,
\[
G_{n+1}(x)\prec (x-E_1)P_n(x).
\]

\medskip \noindent Finally, suppose that
\[
E_1\in I_k,
\qquad
E_2\in I_{k'},
\qquad
k\neq k'.
\] Each of the intervals $I_k$ and $I_{k'}$ already contains one zero
of $F_{n+2}$, namely $E_1$ and $E_2$, respectively. Every other
interior interval contains at least one zero of $P_n$. Hence
$n-2$ zeros of $P_n$ are forced, one in each of the remaining
$n-2$ interior intervals. The two remaining zeros of $P_n$ must
occur together in the exceptional interval. If these two zeros lie in $I_k$ and satisfy
\[
y_{k,n+1}<z_{l,n}<E_1<z_{l+1,n}<y_{k+1,n+1},
\]
then the zero $E_1$ divides $I_k$ into two intervals, each
containing exactly one zero of $P_n$. Every other interval between
consecutive zeros of $(x-E_1)G_{n+1}$ contains exactly one zero of
$(x-E_2)P_n$; in $I_{k'}$, that zero is $E_2$. Therefore, $(x-E_1)G_{n+1}(x)\prec (x-E_2)P_n(x)$, which proves part~(e)(i).

\medskip \noindent Similarly, if the two remaining zeros of $P_n$ lie in $I_{k'}$ and
satisfy
\[
y_{k',n+1}<z_{l,n}<E_2<z_{l+1,n}<y_{k'+1,n+1},
\]
then $(x-E_2)G_{n+1}(x)\prec (x-E_1)P_n(x)$. This proves part~(e)(ii) and completes the proof.
\end{proof}

\end{document}
